\section{The trace coefficient and the passage to all radii}
\label{sec:trace-asymptotic}
Two different normalizations enter the trace asymptotic. Translation Ward
identities fix the amplitude of the stress tensor, while the running Wilson
coefficient fixes the logarithmic scale. After those normalizations, the
numerical coefficient comes from a free curvature contraction. Finally one
must remove the rounded positions and the chosen sequence of scales. We
separate these operations because none can be replaced by choosing a
convenient normalization of the trace after the fact.

\subsection{The free scalar contraction}
Put $N=\dim_{\mathbb R}\mathfrak g$ and use an orthonormal color basis.
The reference potential covariance in four Euclidean dimensions is
$\delta^{AB}\delta_{ab}D(x)$, where
\[
 D(x)=\frac1{4\pi^2|x|^2},\qquad x\neq0.
\]
Indices $a,b,i,j$ range over the four coordinate directions and $A,B$
over the $N$ colors. The curvature covariance obtained by differentiating
this potential is $\delta^{AB}C_{abij}(x)$. Explicitly,
\begin{align*}
 C_{abij}(x)={}&-\delta_{bj}\partial_a\partial_iD(x)
       +\delta_{bi}\partial_a\partial_jD(x)\\
       &+\delta_{aj}\partial_b\partial_iD(x)
       -\delta_{ai}\partial_b\partial_jD(x).
\end{align*}
Write $r=|x|$ and $H_{ai}=\delta_{ai}-2x_ax_i/r^2$. Differentiating
$D$ gives
\[
 \partial_a\partial_iD(x)
   =\frac1{4\pi^2}(-2\delta_{ai}r^{-4}+8x_ax_ir^{-6}),
 \qquad
 C_{abij}(x)=\frac{H_{ai}H_{bj}-H_{aj}H_{bi}}{\pi^2r^4}.
\]
For the second identity expand the products of $H$. Their quartic terms
cancel, leaving exactly the four second-derivative terms above.
Moreover $H=I-2uu^{\mathsf T}$ for the unit vector $u=x/r$, hence
$HH^{\mathsf T}=I$. It follows that
\begin{align*}
 \sum_{i,j}(H_{ai}H_{bj}-H_{aj}H_{bi})^2
 &=2\Bigl(\sum_iH_{ai}^2\Bigr)
       \Bigl(\sum_jH_{bj}^2\Bigr)
       -2\Bigl(\sum_iH_{ai}H_{bi}\Bigr)^2\\
 &=2-2\delta_{ab}.
\end{align*}
Summing over the sixteen ordered pairs $(a,b)$ gives $24$.

The reference scalar is
\[
 O(x)=\frac14\sum_{A,a,b}:F_{ab}^A(x)F_{ab}^A(x):,
\]
where the sum is over ordered coordinate pairs. Here the colons mean the
quadratic Wick field: in its two-point function only contractions between
the two different positions are included. There are exactly two such
pairings. Color contraction gives
$\sum_{A,B}(\delta^{AB})^2=N$, and therefore
\begin{equation}\label{eq:trace-free-scalar}
 \langle O(x)O(0)\rangle_0
 =\frac{2N}{16}\sum_{a,b,i,j}C_{abij}(x)^2
 =\frac{3N}{\pi^4|x|^8}.
\end{equation}
This computation uses the stated reference covariance and Wick rule; the
identification of the constructed reference hierarchy with those data is
an earlier construction requirement. The factor $1/4$ and the use of ordered
pairs are essential to the coefficient.

\subsection{Fixing the physical stress amplitude}
Here is the finite-dimensional step that prevents an arbitrary scalar
multiple of the stress tensor from entering the answer. All vector spaces
in this paragraph are real, including the real spaces of self-adjoint
tensor fields. Let $\mathcal K$ be the space of centered symmetric tensors
with the required rotation symmetry, let $\mathcal V$ be the space in
which their divergences lie, and let
$D:\mathcal K\to\mathcal V$ be the divergence map. The reference theory
has corresponding spaces $\mathcal K_0,\mathcal V_0$ and divergence $D_0$.
The following precise finite-dimensional data are supplied by the stress
comparison estimates: $D_0$ is onto, $\dim\ker D_0=1$,
$\dim\mathcal K=\dim\mathcal K_0$, $\dim\mathcal V=\dim\mathcal V_0$,
and there is a surjective map
$Q:\mathcal K\to\mathcal V_0$ with $\ker Q=\ker D$.
Rank--nullity then gives
\begin{equation}\label{eq:trace-conserved-rank}
 \dim\ker D=\dim\mathcal K-\dim\mathcal V_0
           =\dim\mathcal K_0-\dim\mathcal V_0=1.
\end{equation}
In the construction the surjectivity of $Q$ is obtained from a surjective
finite-dimensional comparison composed with a test map into $\mathcal K$;
a surjective composition is sufficient to make $Q$ itself surjective.
The asserted comparison and kernel equality remain quantitative inputs here.

Let $R\in\ker D$ be the physical centered stress tensor. At scales
$r_i>0$ let $U_i$ be the calibrated isomorphism from reference grade-four
fields to interacting grade-four fields. On tensors set
\[
 Y_i(T)_{ab}=r_i^4U_i^{-1}(T_{ab}).
\]
Suppose the comparison construction supplies eventually conserved tensors
$T_i\in\mathcal K$ and positive numbers $B_i$ with
\begin{equation}\label{eq:trace-trial-stress}
 Y_i(T_i)\longrightarrow R_0,\qquad
 B_i\operatorname{tr}Y_i(T_i)\longrightarrow4d_2O.
\end{equation}
The normalization test is a shell flux functional $\ell_i$ on tensors.
For example, fixing a separated spectator configuration $x$, a smooth
cutoff $\chi$ equal to one near that configuration, a direction $b$, and
one distinguished field $v$, its form is
\[
 \ell_i(T)=-r_i\sum_a\int(\partial_a\chi_i)(z)
   S_{T_{ab},v_i,\boldsymbol q_i}(z,r_ix)\,dz,
 \qquad \chi_i(z)=\chi(z/r_i).
\]
The derivatives of $\chi_i$ are supported away from the spectators, so the
integrals involve separated correlations. The translation Ward identity
for the physical tensor identifies this with a derivative insertion.
The needed convergence of that insertion and of the comparison shell
integral has the following exact consequence:
\begin{equation}\label{eq:trace-flux-test}
 \ell_i(R)\longrightarrow y\neq0,
 \qquad \ell_i(T_i)\longrightarrow y.
\end{equation}
The same nonzero $y$ occurs because $R_0$ has the reference translation
Ward identity. The existence of a detecting reference derivative field,
and the jet estimates justifying these two limits, are part of the stress
comparison input, rather than consequences of rank--nullity.

\begin{lemma}\label{lem:trace-normalized-stress}
Equations \eqref{eq:trace-conserved-rank}--\eqref{eq:trace-flux-test} imply
\[
 Y_i(R)\longrightarrow R_0,\qquad
 B_i\operatorname{tr}Y_i(R)\longrightarrow4d_2O.
\]
\end{lemma}
\begin{proof}
The first limit in \eqref{eq:trace-flux-test} implies $R\neq0$.
Since the conserved space is one-dimensional, eventually $T_i=c_iR$
for a real scalar $c_i$. Linearity of $\ell_i$ gives
$c_i\ell_i(R)=\ell_i(T_i)$. For large $i$ the denominator $\ell_i(R)$
is nonzero, and the two limits give $c_i\to y/y=1$. In particular $c_i$
is eventually nonzero and $c_i^{-1}\to1$. Multiply both limits in
\eqref{eq:trace-trial-stress} by $c_i^{-1}$ and use linearity of $Y_i$
and of the trace. This yields exactly the asserted limits.
\end{proof}

Centering the trace in the chosen phase can differ from tracing the
mixed-state-centered tensor $R$ by a fixed multiple $c\mathbf1$ of the
unit. If $U_i\mathbf1_0=\mathbf1$ eventually and $B_ir_i^4\to0$, then
\[
 U_i^{-1}(B_ir_i^4c\mathbf1)=B_ir_i^4c\mathbf1_0\longrightarrow0.
\]
Consequently the selected centered trace $\Theta$ has the same limit
\begin{equation}\label{eq:trace-moving-label}
 a_i:=U_i^{-1}(B_ir_i^4\Theta)\longrightarrow4d_2O.
\end{equation}
The identities above explain how the physical Ward normalization, rather
than a later rescaling of $\Theta$, determines this coefficient.

\subsection{Rounded correlations and moving labels}
We next state exactly the remaining approximation data used to pass from
\eqref{eq:trace-moving-label} to the trace asymptotic. Write
$S_\chi(v,w;x_0,x_1)$ for the selected two-point function and
$S_0(v,w;x_0,x_1)$ for the reference one. Both are bilinear in the field
labels. The selected centered trace is a scalar under Euclidean rotations,
and translations preserve its correlations. Thus its two-point function
has the form
\begin{equation}\label{eq:trace-radial}
 S_\chi(\Theta,\Theta;x_0,x_1)=T(|x_0-x_1|)
       \quad(x_0\neq x_1).
\end{equation}
To justify this reduction, an orthogonal map sends the nonzero difference
$x_0-x_1$ to $-|x_0-x_1|e_0$; equal-length vectors are related by a
reflection, with the identity covering equality. A translation then sends
the two points to $-te_0/2,te_0/2$, where $t=|x_0-x_1|$.
Covariance and the scalar transformation law of $\Theta$ give
\eqref{eq:trace-radial}.

\begin{assumption}[Trace comparison along scales]\label{input:trace-scales}
For every sequence $s_n>0$ tending to zero there is a subsequence
$r_i=s_{n_i}$, with $n_i$ strictly increasing, and comparison data with
the following properties. The maps $U_i$ and numbers $B_i>0$ satisfy
\eqref{eq:trace-moving-label} and
\[
 \frac{\log(r_i^{-1})}{B_i}\longrightarrow\frac1{2b_0},\qquad
 b_0=\frac{11C_{\mathfrak g}}{48\pi^2}>0,\qquad d_2=-\frac{b_0}{2},
\]
where $C_{\mathfrak g}>0$ is the adjoint Casimir in the fixed normalization.
There are rounded configurations $q_i(x)\to x$ and errors
$\varepsilon_i\to0$. Let $E_0$ be the reference grade-four field space
and fix its finite basis $(e_\alpha)$. At every coordinatewise separated
configuration $x$ in the interior of the comparison window, eventually
\begin{align}
 &\bigl|S_\chi(U_ie_\alpha,U_ie_\beta;r_iq_i(x))
                    -S_0(e_\alpha,e_\beta;q_i(x))\bigr|\nonumber\\
 &\hspace{30mm}\leq\varepsilon_i\bigl(1+M_{\alpha\beta}(q_i(x))\bigr),
 \label{eq:trace-basis-error}
\end{align}
where $M_{\alpha\beta}$ is continuous near $x$. The reference correlations
are continuous there. The comparison window includes a neighborhood of
$x(t)=(0,(t/2,t/2,t/2,t/2))$ for each $1/2<t<2$.
\end{assumption}
This input is the ultraviolet and stress comparison problem. The positive
trace spectral measure and its Laplace formula are supplied by
Section~\ref{sec:trace-spectrum} for the same reconstructed centered trace.
We have not inferred the comparison estimates from compactness.

\begin{lemma}\label{lem:trace-moving-correlation}
Under Input~\ref{input:trace-scales}, for each eligible fixed $x$,
\[
 r_i^8\log^2(r_i^{-1})T(r_i|q_i(x)_0-q_i(x)_1|)
       \longrightarrow \frac{A}{|x_0-x_1|^8},
 \qquad A=\frac{3N}{\pi^4}.
\]
\end{lemma}
\begin{proof}
The right side of \eqref{eq:trace-basis-error} tends to zero, since
$M_{\alpha\beta}(q_i(x))$ converges and is therefore bounded.
Continuity of $S_0$ gives convergence of the left correlation to
$S_0(e_\alpha,e_\beta;x)$. This also holds for convergent, varying labels:
expand each of them in the fixed finite basis. Every coordinate converges,
and bilinearity expresses the correlation as a finite sum of products of
two coordinates and one of the convergent basis correlations. In particular,
\eqref{eq:trace-moving-label} gives
\[
 S_\chi(U_ia_i,U_ia_i;r_iq_i(x))
 \longrightarrow (4d_2)^2 S_0(O,O;x).
\]
The definition of $a_i$, bilinearity and \eqref{eq:trace-radial} identify
the expression on the left with
$B_i^2r_i^8T(r_i|q_i(x)_0-q_i(x)_1|)$.
Multiply by $(\log(r_i^{-1})/B_i)^2$ and use
\eqref{eq:trace-free-scalar}. The limiting coefficient is
\[
 \frac{(4d_2)^2}{(2b_0)^2}\frac{3N}{\pi^4}
       =\frac{12d_2^2N}{b_0^2\pi^4}
       =\frac{3N}{\pi^4}.
\]
All divisions are legitimate since $b_0>0$ and $B_i>0$.
\end{proof}

\subsection{Removing rounding and subsequences}
The rounded separation need not equal the intended radius. Pointwise
convergence at rounded arguments alone does not identify $T(r_i)$.
The positive spectral representation resolves this: $T$ is real,
nonnegative and decreasing, because $e^{-tp_0}$ has those properties for
each $p_0\geq0$. Integrability at the smaller positive time permits the
comparison under the integral.

\begin{theorem}\label{thm:trace-asymptotic}
Input~\ref{input:trace-scales}, with the reference normalization above,
implies
\[
 \lim_{r\downarrow0}r^8\log^2(r^{-1})T(r)=\frac{3N}{\pi^4}>0.
\]
\end{theorem}
\begin{proof}
First fix any one comparison subsequence $r_i$. Let
$c_i=r_i^8\log^2(r_i^{-1})\geq0$. For $1/2<t<2$ take the configuration
$x(t)$ in the input. Its Euclidean separation is exactly $t$, all four
coordinate separations are nonzero, and
$w_i(t):=|q_i(x(t))_0-q_i(x(t))_1|\to t$.
Lemma~\ref{lem:trace-moving-correlation} gives
\[
 c_iT(r_iw_i(t))\longrightarrow A/t^8.
\]
If $t>1$, eventually $w_i(t)>1$, so monotonicity yields
$c_iT(r_i)\geq c_iT(r_iw_i(t))$. If $t<1$, eventually
$0<w_i(t)<1$ and the reverse inequality holds. Given $\epsilon>0$,
choose $t_+>1$ and $t_-<1$ sufficiently close to one that
$|A/t_\pm^8-A|<\epsilon/2$. For all sufficiently large $i$, the two
comparison values are within $\epsilon/2$ of their limits. They bound
$c_iT(r_i)$ between $A-\epsilon$ and $A+\epsilon$. Thus
$c_iT(r_i)\to A$.

To obtain the limit at all positive radii, suppose it failed. There would
be $\epsilon>0$ and a sequence $s_n\in(0,1/n)$, $n\geq1$, for which
$|s_n^8\log^2(s_n^{-1})T(s_n)-A|\geq\epsilon$ for every $n$.
Input~\ref{input:trace-scales} supplies a comparison subsequence of this
same sequence. The preceding argument makes its weighted values converge
to $A$, contradicting that inequality. Finally $N>0$ for the nonabelian
simple Lie algebra under consideration, so $A>0$.
\end{proof}
Combining this theorem with Theorem~\ref{thm:profile-exclusion} proves the
finite-profile exclusion from the same trace comparison input. The rank,
Ward, approximation and scale-selection estimates identified above still
need to be supplied by the construction; the algebraic and limiting
passages do not replace them.
